\documentclass[11pt]{article}
\usepackage[margin=1in]{geometry}
\usepackage{amsmath, amssymb, amsthm, bm, mathtools}
\usepackage{booktabs}
\usepackage{array}
\usepackage{enumitem}

\title{\textbf{Quiz 2} \vspace{ 0.4cm }\\ECON 101: Macroeconomics }
\author{Instructor: Muhebullah Karimzada}
\date{February 05,  2026}

\begin{document}
\maketitle



\section*{Instructions}
Answer all questions on the answer sheet.  Show all work for numerical problems. Multiple-choice questions have one correct answer.\\



\textbf{1.} Suppose beef and chicken are substitutes. If the price of beef rises, holding all else constant:  
\begin{enumerate}[label=\alph*)]
    \item Demand for chicken decreases.  
    \item Supply of chicken decreases.  
    \item Demand for chicken increases.  
    \item Quantity demanded of chicken decreases.  
\end{enumerate}

\bigskip

\textbf{2.} Which of the following transactions is \textbf{included} in Canada’s GDP?  
\begin{enumerate}[label=\alph*)]
    \item A Canadian student buys used textbooks from another student.  
    \item A Canadian firm sells lumber to a U.S. company.  
    \item A Canadian tourist purchases a meal in Paris.  
    \item A Canadian household does its own home repairs.  
\end{enumerate}

\bigskip

\textbf{3.} Real GDP differs from nominal GDP because real GDP:  
\begin{enumerate}[label=\alph*)]
    \item Includes only goods, not services.  
    \item Is adjusted for population size.  
    \item Values output at constant prices.  
    \item Excludes investment spending.  
\end{enumerate}

\bigskip
\newpage
\textbf{4.} If the GDP deflator increases from 120 to 132, the inflation rate is approximately:  
\begin{enumerate}[label=\alph*)]
    \item 10\%  
    \item 12\%  
    \item 9\%  
    \item 8\%  
\end{enumerate}

\bigskip

\textbf{5.} Which of the following would \textbf{not} increase GDP?  
\begin{enumerate}[label=\alph*)]
    \item Higher consumption of new cars.  
    \item More exports of manufactured goods.  
    \item An increase in unpaid household work.  
    \item Greater investment in factories.  
\end{enumerate}

\bigskip



\textbf{Problem 1:}  
The market demand and supply for pizzas in Toronto are given by:  
\[
Q_D = 200 - 5P, \quad Q_S = 20 + 3P,
\]  
where \(Q\) is the number of pizzas and \(P\) is the price in dollars.  

\begin{enumerate}[label=\alph*)]
    \item Find the equilibrium price and quantity.  
    \item If the government imposes a price floor at \$25, calculate the resulting surplus or shortage.  
\end{enumerate}

\bigskip

\textbf{Problem 2:}  
Suppose the following Canadian macroeconomic data (in billions):  

\begin{itemize}
    \item Consumption: \$1,200  
    \item Investment: \$500  
    \item Government spending: \$600  
    \item Exports: \$450  
    \item Imports: \$550  
\end{itemize}

\begin{enumerate}[label=\alph*)]
    \item Calculate nominal GDP using the expenditure approach.  
    \item If population = 37 million, compute GDP per capita.  
    \item If potential GDP = \$3,000 billion, calculate the output gap as a percentage.  
\end{enumerate}


\end{document}
